\section{Inner Structure and Matter}
\label{sec:matter}

This is the frontier of the model. A scalar ternary tone gives a vibe a state but not the inner structure that matter needs, namely spin, chirality, fermions, antiparticles, charge, and the exclusion principle. This section sketches how that inner structure can be built without adding any new primitive. The constructions are grounded in known mathematics but are not yet proven on a hyperbolic mesh, so this section is openly speculative.

\subsection{The monism requirement on the inside}

Matter cannot be added as a new primitive. By monism (Section~\ref{sec:constraints}) there is nothing but vibes, tones, notes, and beats, so any internal degree of freedom, such as spin or charge, must be a derived pattern of a small system of vibes and notes, not a substance inside one vibe. ``A vibe carries spin'' is shorthand for ``a small patch of vibes and notes has a mode that behaves like spin.''

\subsection{The spinor as a graded complex of vibes}

The cleanest construction is the Kähler-Dirac form, in which a spinor is not a mysterious column of components but a tone on each cell of a small patch, read here as a graded complex of vibes. The vertices are base vibes. The edges are relational vibes, the notes themselves. The faces are higher relational vibes over small loops. A spinor is then a tuple of vibe tones across these levels on one neighborhood, knit together by the discrete exterior derivative, which is just the note-adjacency between a relational vibe and the vibes it relates. Every component is a vibe. The number of components is fixed by the cell combinatorics of one tile-neighborhood, not chosen by hand, which is also how the model would address fermion doubling. The gamma matrices that define a spinor become the operation of raising or lowering the relational order, binding or releasing a relational vibe, both pure note operations.

\begin{proposition}[The belt trick]
Hold one end of a belt fixed and rotate the buckle one full turn. The belt is twisted and cannot be undone without rotating back. Rotate the buckle two full turns and the twist slides off without rotating the buckle at all. The buckle's ``spin,'' its $720$-degree periodicity, lives entirely in its connection to its surroundings, not in the buckle itself. A spinor in the mesh is the same. The double-cover sign is a holonomy of the notes tethering a patch to the rest, not a hidden coin inside one vibe.
\end{proposition}

The complex phase a spinor uses is not primitive. By discreteness it is carried by a finite clock-tone, a $\mathbb{Z}_q$ value on a relational vibe, or by the parity of the beat, advancing one step per beat. This same internal tick is a candidate for mass, the rest-frame oscillation rate. Full exchange statistics, the minus sign that makes matter rigid, would come from braiding two such clusters, which is natural on the two-dimensional hyperbolic mesh.

\subsection{Toward charge}

The same scaffold extends. An internal charge, such as electric or color charge, is a finite representation carried by the vibe-structure, transforming under a finite internal symmetry of the rule, and its conserved value is the most literal karma the mesh has. Discreteness predicts that the allowed charges should be a small finite set rather than a continuum, the right shape for the particle zoo. A particle, in the end, may be best read as a stable topological pattern, a knot in the tone field, whose inner structure is a topological invariant and whose stability is protected by a winding number, rather than a special fundamental vibe.
